\begin{statement}[C.8.12; Proved: a radial derivative identity]
Put $p=\wp'_\tau(\mu)$ and
\[
 X=\frac{3Wp}{g_2},\qquad F=1-X,
 \qquad W=\frac{\vartheta_1'(\mu,\tau)}{\vartheta_1(\mu,\tau)}
              +\frac{2\pi i\Im\mu}{\Im\tau}.
\]
These are single-valued modular invariants wherever $g_2\ne0$.
They have continuous extensions at $t=0,1$, with
$X(0)=1/4$ and $X(1)=0$. At every $t\ne0,1$,
\begin{equation}\label{eq:cc-radial}
 (t-1)L_t=\frac{X(X-2)}{24},\qquad
 r\frac{d}{dr}L(1+re^{i\theta})
       =\frac{\Re(X(X-2))}{12}.
\end{equation}
Also, if $a=\Im\mu/\Im\tau$ and $R=\mu'-a$, then
\begin{equation}\label{eq:cc-FR}
 F=\frac{6\pi ip}{g_2}R.
\end{equation}
Consequently $|F|<1$ implies strict decrease along the ray
from $1$ through $t$.
\end{statement}
\begin{proof}
A lattice translation leaves $W$ unchanged, and replacing $\mu$
by $-\mu$ changes the signs of both $W$ and $p$. Under a modular
change, $W,p,g_2$ have respective weights $1,3,4$, so $X$ is
invariant. The two identities of C.8.11 give
\[
 L_\tau=\frac{iW(3Wp-2g_2)}{24\pi p},\qquad
 R=\frac{g_2-3Wp}{6\pi ip}.
\]
Ramanujan's identities and $p^2=-g_3$ give
\[
 t'_\tau=-2\pi it\frac{E_6}{E_4},\qquad
 \frac{t-1}{t'_\tau}=-\frac{E_6}{2\pi iE_4}.
\]
Substitution, using $g_2=4\pi^4E_4/3$ and
$g_3=8\pi^6E_6/27$, proves \eqref{eq:cc-radial} and
\eqref{eq:cc-FR}. Since $X(X-2)=F^2-1$, $|F|<1$ makes
its real part strictly negative.

At $i$, $W=p=0$ and $g_2\ne0$, so $X\to0$. At $\rho$, put $\delta=\tau-\rho$ and $p_0=\wp'_\rho(\mu_0)$.
The expansion in C.3.2 and Ramanujan's identities give
\[
 R=\frac{3p_0}{8\pi^2}\delta+O(|\delta|^2),\qquad
 g_2=\frac{3ip_0^2}{\pi}\delta+O(\delta^2).
\]
Substitution into \eqref{eq:cc-FR} gives $F(\rho)=3/4$.
Equivalently, the two possible limits furnished by
\eqref{eq:cc-radial} and $L_t(0)=7/384$ are $X=1/4,7/4$;
the expansion of $R$ selects $1/4$. The same expansions show
continuity on the coarse $t$-line.
\end{proof}
